\section{Debit-Budgeted CFR}
\label{sec:method}

DB-CFR separates execution, settlement, and planning.  The \emph{executor}
produces an event-faithful learner record for one declared width and order.
The \emph{auditor} computes the target debit of an executed partial block.
The \emph{planner} reads only completed learner history and settled charges;
it never evaluates an unselected future profile.  Complete coverage connects
these modules by supplying an always-available zero-debit action.

\begin{table}[h]
\caption{One causal DB-CFR block.  Partial plans execute only while the settled
account is solvent and the audit queue has capacity.}
\label{tab:db-cfr-algorithm}
\centering
\small
\begin{tabularx}{\linewidth}{@{}p{0.055\linewidth}X@{}}
\toprule
\multicolumn{2}{@{}l@{}}{\textbf{Input:} learner state, settled account
$A_{k-1}$, pending set $P_{k-1}$, budget $b(T_k)$, cap $\ell$.} \\
\midrule
1 & Settle every completed audit, updating $A_{k-1}$ and $P_{k-1}$. \\
2 & If $A_{k-1}>b(T_k)$ or $|P_{k-1}|\ge\ell$, choose complete coverage. \\
3 & Otherwise choose one registered plan by a fixed completed-history rule. \\
4 & Execute that plan once under the event-faithful contract. \\
5 & Settle a complete block at charge zero; enqueue a partial block for
target-debit audit. \\
\bottomrule
\end{tabularx}
\end{table}

The method is causal because Step 3 cannot read the current block's charge.
Theorem~\ref{thm:debit-controller} permits any such history rule: cost,
previous charges, and learner state may affect the plan, while the debit
account controls when exact recovery is required.  A final reported
certificate either drains the queue or retains the reserve
$R_K^{(\ell)}$.

\subsection{Computing the settled charge}

The charge is defined by the dense target at the learner origins that actually
occurred.  An exact audit evaluates the omitted components at these immutable
origins and applies Equation~\ref{eq:block-charge}; interval or factored
observers may instead upper-bound the same terms.  The appendix gives a finite
summation-by-parts identity that separates the weighted endpoint from component
motion along the executed path.  Complete coverage uses the stronger eventwise
identity $h_{\rho,t}=g_{\rho,t}$ and therefore settles at zero without an
audit.

\subsection{Lagged causal instantiation}

The public-Leduc implementation uses synchronous post-block settlement
($\ell=1$), widths $2$, $5$, and $10$, twelve ten-round blocks, and
$b(T)=12\sqrt T$.  When the account is insolvent it executes width 10.  While
solvent, it observes widths 5 and 2 once, then chooses the least-work partial
width whose most recently settled charge fits the current slack; if neither
does, it chooses the width with the smaller last charge.  The selected block
executes once and its exact charge settles before the next decision.

This rule receives neither exploitability nor a future candidate state.  It
uses zero unselected candidate replays, and its selected width can change with
the realized learner path.  The appendix retains a prospective exact-lookahead
controller as a strong comparator and audits its cost-ordered offer search.

\subsection{Bounded asynchronous settlement}

When selected-path audits run concurrently, the learner may open up to
$\ell$ partial blocks before their charges settle.  Each audit consumes an
immutable event record; the planner enters complete recovery when either the
settled account or the queue cap blocks another partial action.  The finite
reserve in Theorem~\ref{thm:debit-controller} accounts for every hidden charge,
so asynchronous completion changes the bound but not the declared learner
recurrence.  This interface removes audit latency from the sequential decision
path; audit CPU and record traffic remain part of the system cost measured in
Section~\ref{sec:experiments}.
